\textbf{Classroom test, part 1.} You have 40 minutes in total.
Two wide-field maps, around $100^\circ$ wide, are provided. Map 1 is for an observer of the southern
hemisphere and Map 2 for an observer of the northern hemisphere. The Milky Way is not depicted.
Write your answers on the maps themselves. You can divide the available time between the two maps as you wish.

If more than the requested number of answers is provided, the following rules apply.
For questions 1a, 1b, 2a and 2b there is no penalty for one extra answer, that is, marking $N+1$ objects when
$N$ are requested. Each extra object after that loses partial points. Wrong answers do not receive negative marking.
Whenever drawing constellations, use any common stick figure from atlases or planetarium software.
All reasonable drawings are accepted, but keep them simple and recognisable. For example, Andromeda can
be represented as the two back legs of Pegasus, rather than a person as in the IAU maps.
The source gives the following two acceptable drawings of Virgo.
\ioaafig{2023_OBS-SKY_Q13-f1.pdf}
\textbf{Map 1, southern sky.}
\begin{description}
\item[(a)] Draw stick figures for any 10 constellations on the map and write their IAU code next to them. \hfill\mbox{[10 points]}
If you do not know the stick figure but know its area, you may provide its IAU code there.
Refer to the supplied IAU list. Only choose constellations that are almost completely visible.
Partially cut-off constellations are not accepted.
\item[(b)] Identify any 20 stars and provide their proper names or Bayer designations in English. \hfill\mbox{[10 points]}
Both names are accepted but knowing both gives no extra points. The stars need not be in the constellations chosen in (a).
\item[(c)] A very bright star is missing. Mark its position with an ``x''. \hfill\mbox{[2.5 points]}
\end{description}
\textbf{Map 2, northern sky.}
\begin{description}
\item[(a)] Draw stick figures of any 10 constellations and write their IAU codes next to them, using the same instructions as 1a. \hfill\mbox{[10 points]}
\item[(b)] Identify 20 stars and provide their proper names or Bayer abbreviations in English, using the same instructions as 1b. \hfill\mbox{[10 points]}
\item[(c)] A supernova is present. Identify and circle it. \hfill\mbox{[2.5 points]}
\end{description}
The IAU constellation list and the two maps follow.
\ioaafig{2023_OBS-SKY_Q13-f2.pdf}\ioaafig{2023_OBS-SKY_Q13-f3.pdf}\ioaafig{2023_OBS-SKY_Q13-f4.pdf}
